\subsection{主不等式的最简等价形式：$\tau=\beta-(1-\alpha)$}\label{sec:r17:tau}

第~\ref{sec:harm} 节把问题写成 $\rho=\alpha+\beta\le\rho_0$，其中 $\alpha=x^\top Qx=|m|^2\le1$、
$\beta=x^\top Q_-x\ge0$。本轮把它进一步压缩：\textbf{需要控制的其实只是一个标量差}。

\begin{lemma}[$\tau$-形式]\label{lem:tauform}
记 $\tau:=\rho-1$、$\tau_0:=\rho_0-1=\frac{\sqrt5+2\sqrt{35}-9}{18}=0.281568196872168$。则
\begin{equation}\label{eq:tauform}
Q_-x=\tau x+(I-Q)x,\qquad\text{从而}\qquad
\boxed{\ \tau=\beta-(1-\alpha)=x^\top Q_-x-\|(I-Q)x\|^2\ } .
\end{equation}
于是\textbf{主猜想 $\kappa_3\le\rho_0$ 等价于：对一切 $N$ 与一切 Perron 对 $(Q_+,x)$，}
\begin{equation}\label{eq:maincrux}
x^\top Q_-x-\|(I-Q)x\|^2\ \le\ \tau_0 .
\end{equation}
\end{lemma}

\begin{proof}
$Q_+=Q+Q_-$ 且 $Q_+x=\rho x$ 给 $Q_-x=\rho x-Qx=(\rho-1)x+(I-Q)x$。
与 $x$ 作内积：$\beta=\tau+\|(I-Q)x\|^2$（用 $(I-Q)x\perp x$），
而 $\|(I-Q)x\|^2=x^\top(I-Q)x=1-\alpha$（$Q$ 是秩三正交投影，$Q^2=Q$）。
\end{proof}

\begin{remark}[缺口在哪里]\label{rmk:taugap}
\eqref{eq:maincrux} 的左端是两项之差：
\emph{钝部能量} $x^\top Q_-x=\beta$ 减去 \emph{Perron 向量离开列空间的距离} $\|(I-Q)x\|^2=1-\alpha$。
冠军处 $\beta=0.304972937817$、$1-\alpha=0.023404740945$，差恰为 $0.281568196872=\tau_0$
（数值核对到 $1.22\times10^{-15}$）——\textbf{两项都不小，是它们的差在取等}。
这解释了为什么过去所有「单独控制 $\beta$」的尝试（例如 $\beta\le\lambda_{\max}(Q_-)$）都注定松弛：
冠军处 $\lambda_{\max}(Q_-)=0.379920437>\tau_0$，
单独控制 $\beta$ 最好也只能给 $\rho\le1.37992$，比 Mohar 的 $1.36603$ 还差。
\textbf{必须让 $(I-Q)x$ 这一项参与抵消。}
\end{remark}

\begin{corollary}[$1+\lambda_{\max}(Q_-)$ 界的一行新证]\label{cor:oneneg}
$\rho\le1+\lambda_{\max}(Q_-)$。
\end{corollary}

\begin{proof}
由 \eqref{eq:tauform}，$\tau=\beta-\|(I-Q)x\|^2\le\beta=x^\top Q_-x\le\lambda_{\max}(Q_-)$。
\end{proof}

\begin{lemma}[$Q_-x$ 的精确范数]\label{lem:normqmx}
$\|Q_-x\|^2=(1-\alpha)(1+2\tau)+\tau^2=\beta^2+\alpha(1-\alpha)$，
且 $(I-Q)Q_-x=\rho(I-Q)x$、$QQ_-x=\tau Qx$。
\end{lemma}

\begin{proof}
第一式：把 \eqref{eq:tauform} 取范数平方，用 $(I-Q)x\perp x$。
第二式：$(I-Q)$ 作用在 \eqref{eq:tauform} 两端（$(I-Q)Q=0$）。
第三式：$Q$ 作用在两端（$Q(I-Q)=0$、$Qx=Qx$）。
末一等式 $\|Q_-x\|^2=\beta^2+\alpha(1-\alpha)$ 就是第 15 轮的 $\sigma$-恒等式
$\sum_ig_i\sigma_i^2=\beta^2+\alpha(1-\alpha)$（因 $(Q_-x)_i=\sqrt{g_i}\,\sigma_i$），
两者在此互相印证。
\end{proof}

\subsection{精确 Frobenius 恒等式与一行式谱地板}\label{sec:r17:frob}

\begin{lemma}[正部与负部的能量守恒]\label{lem:frob}
记 $e:=\|Q_-\|_F^2$。则
\begin{equation}\label{eq:frob}
\langle Q,Q_-\rangle=-e,\qquad \|Q_+\|_F^2=3-e,\qquad \sum_k\lambda_k(Q_+)^2=3-e,\qquad \sum_k\lambda_k(Q_+)=3 .
\end{equation}
\end{lemma}

\begin{proof}
$(Q_-)_{ij}>0$ 仅当 $Q_{ij}<0$，此时 $Q_{ij}=-(Q_-)_{ij}$，故
$\langle Q,Q_-\rangle=\sum_{ij}Q_{ij}(Q_-)_{ij}=-\sum(Q_-)_{ij}^2=-e$（\emph{精确}等式，非数值）。
$Q_+=Q+Q_-$ 给 $\|Q_+\|_F^2=\|Q\|_F^2+2\langle Q,Q_-\rangle+\|Q_-\|_F^2=3-2e+e=3-e$，
其中 $\|Q\|_F^2=\tr Q^2=\tr Q=3$（$Q$ 是秩三正交投影，即 Welch 界取等）。
迹：$\tr Q_+=\tfrac12(\tr Q+\tr|Q|)=\tfrac12(3+3)=3$，因 $\tr|Q|=\sum|Q_{ii}|=\sum g_i=3$。
\end{proof}

\begin{remark}[这条恒等式的价值]\label{rmk:frobval}
\eqref{eq:frob} 把「正部的谱」与「负部的能量」用一个\textbf{等式}锁在一起：
$\rho$ 想大，$e$ 就必须小；$e$ 小，则 $\|Q_+\|_F^2=3-e$ 大，正部的谱就必须更「铺开」。
冠军处 $e=\rho_0\tau_0=\frac{16+5\sqrt7}{81}=0.360848846362012$，
$\|Q_+\|_F^2=2.639151153638$，与 $\sum\lambda_k^2$ 直接计算吻合到 $1.3\times10^{-15}$。
\end{remark}

\begin{corollary}[一行式谱地板]\label{cor:floor2}
$n_+(Q_+)\ge\dfrac{9}{3-e}$（$n_+$ $=$ 正特征值个数）。
特别地，若 $e\le\rho_0\tau_0$ 则 $n_+(Q_+)\ge4$。
\end{corollary}

\begin{proof}
$3-e=\sum\lambda_k^2\ge(\sum\lambda_k)^2/n_+=9/n_+$。第二个断言：$9/(3-0.360848846)=3.410\Rightarrow n_+\ge4$。
\end{proof}

\begin{remark}
这与第 15 轮的谱地板 $n_+\ge1+(3-\rho)^2/(3-\rho^2)$（在 $\rho=\rho_0$ 处同为 $4$）结论一致，
但证明只有一行，且\textbf{不需要} $\rho$ 的取值，只需要 $e$。两条地板互为对偶。
\end{remark}

\begin{proposition}[条件性重排上界，优于 \STW]\label{prop:majbound}
若 $Q_+\succeq0$ 且 $n_+(Q_+)\le4$，则
\begin{equation}\label{eq:majbound}
\rho\ \le\ \frac{3+\sqrt{9-12e}}{4}\ \le\ \frac34+\frac{\sqrt{9-12\cdot0}}4=\frac32 ,
\end{equation}
且在冠军的 $e=\rho_0\tau_0$ 处右端 $=1.290244$，\textbf{严格小于 \STW\ 的 $\Phi=(3+\sqrt5)/4=1.309017$}。
\end{proposition}

\begin{proof}
设 $Q_+$ 的正特征值为 $\rho=\lambda_1>\lambda_2,\lambda_3,\lambda_4\ge0$（其余为 $0$）。
则 $\sum_{k=2}^4\lambda_k=3-\rho$、$\sum_{k=2}^4\lambda_k^2=3-e-\rho^2$，
Cauchy--Schwarz 给 $3-e-\rho^2\ge(3-\rho)^2/3$，整理得 $4\rho^2-6\rho+3e\le0$，解出 \eqref{eq:majbound}。
\end{proof}

\begin{remark}[诚实标注]\label{rmk:majcond}
命题~\ref{prop:majbound} 是\textbf{条件性}的：冠军自己 $n_+=6$，不满足 $n_+\le4$，
所以 \eqref{eq:majbound} 在冠军处\textbf{不适用}（把 $n_+=6$ 代入同样的重排论证只得 $1.474$）。
它的意义是：\textbf{如果}能把 $n_+$ 压到 $4$（推论~\ref{cor:floor2} 说 $4$ 是下界，所以 $n_+=4$ 是唯一可能的极值层），
那么 $\rho\le1.290244$ 就已经\textbf{无条件地优于} \STW。这把「改进 \STW」与「压低 $n_+$」两件事
精确地联系起来了，也给出了下一轮的一条明确路线（见 \S\ref{sec:r1517:status} 的 P2）。
\end{remark}

